\documentclass[10pt,a4paper]{amsart}
\usepackage{amsmath,amssymb,amsthm,graphicx}
\usepackage{verbatim}
\newenvironment{reference}{\comment}{\endcomment}
%\newenvironment{reference}{}{}
\newenvironment{slogan}{\comment}{\endcomment}
\newenvironment{history}{\comment}{\endcomment}

% For commutative diagrams we use Xy-pic
\usepackage[all]{xy}

% We use 2cell for 2-commutative diagrams.
\xyoption{2cell}
% Theorem environments.
%
\theoremstyle{plain}
\newtheorem{theorem}[subsection]{Theorem}
\newtheorem{proposition}[subsection]{Proposition}
\newtheorem{lemma}[subsection]{Lemma}

\theoremstyle{definition}
\newtheorem{definition}[subsection]{Definition}
\newtheorem{example}[subsection]{Example}
\newtheorem{exercise}[subsection]{Exercise}
\newtheorem{situation}[subsection]{Situation}

\theoremstyle{remark}
\newtheorem{remark}[subsection]{Remark}
\newtheorem{remarks}[subsection]{Remarks}

\numberwithin{equation}{subsection}

% Macros
%
\def\lim{\mathop{\mathrm{lim}}\nolimits}
\def\colim{\mathop{\mathrm{colim}}\nolimits}
\def\Spec{\mathop{\mathrm{Spec}}}
\def\Hom{\mathop{\mathrm{Hom}}\nolimits}
\def\Ext{\mathop{\mathrm{Ext}}\nolimits}
\def\SheafHom{\mathop{\mathcal{H}\!\mathit{om}}\nolimits}
\def\SheafExt{\mathop{\mathcal{E}\!\mathit{xt}}\nolimits}
\def\Sch{\mathit{Sch}}
\def\Mor{\mathop{\mathrm{Mor}}\nolimits}
\def\Ob{\mathop{\mathrm{Ob}}\nolimits}
\def\Sh{\mathop{\mathit{Sh}}\nolimits}
\def\NL{\mathop{N\!L}\nolimits}
\def\CH{\mathop{\mathrm{CH}}\nolimits}
\def\proetale{{pro\text{-}\acute{e}tale}}
\def\etale{{\acute{e}tale}}
\def\QCoh{\mathit{QCoh}}
\def\Ker{\mathop{\mathrm{Ker}}}
\def\Im{\mathop{\mathrm{Im}}}
\def\Coker{\mathop{\mathrm{Coker}}}
\def\Coim{\mathop{\mathrm{Coim}}}

% Boxtimes
%
\DeclareMathSymbol{\boxtimes}{\mathbin}{AMSa}{"02}

%
% Macros for moduli stacks/spaces
%
\def\QCohstack{\mathcal{QC}\!\mathit{oh}}
\def\Cohstack{\mathcal{C}\!\mathit{oh}}
\def\Spacesstack{\mathcal{S}\!\mathit{paces}}
\def\Quotfunctor{\mathrm{Quot}}
\def\Hilbfunctor{\mathrm{Hilb}}
\def\Curvesstack{\mathcal{C}\!\mathit{urves}}
\def\Polarizedstack{\mathcal{P}\!\mathit{olarized}}
\def\Complexesstack{\mathcal{C}\!\mathit{omplexes}}
% \Pic is the operator that assigns to X its picard group, usage \Pic(X)
% \Picardstack_{X/B} denotes the Picard stack of X over B
% \Picardfunctor_{X/B} denotes the Picard functor of X over B
\def\Pic{\mathop{\mathrm{Pic}}\nolimits}
\def\Picardstack{\mathcal{P}\!\mathit{ic}}
\def\Picardfunctor{\mathrm{Pic}}
\def\Deformationcategory{\mathcal{D}\!\mathit{ef}}
\usepackage{cleveref}
\crefname{originallemma}{original lemma}{original lemmas}
\Crefname{originallemma}{Original lemma}{Original lemmas}
\makeatletter
\newlabel{lemma-tensor-finiteness}{{05BS}{source}{}{sourcepointer-1}{}}
\newlabel{lemma-tensor-finiteness@cref}{{[originallemma][0][]05BS}{[1][0][]source}{}{}{}}
\makeatother
\crefname{originallemma}{original lemma}{original lemmas}
\Crefname{originallemma}{Original lemma}{Original lemmas}
\makeatletter
\newlabel{proposition-ML-characterization}{{059E}{source}{}{sourcepointer-2}{}}
\newlabel{proposition-ML-characterization@cref}{{[originallemma][0][]059E}{[1][0][]source}{}{}{}}
\makeatother
\title{The tensor product of two Mittag-Leffler modules is Mittag-Leffler}
\author{The Stacks Project Authors}
\date{stacks-project:algebra.tex:local-snapshot-20261012; GNU FDL 1.2 or later}
\begin{document}
\maketitle
\noindent\textit{Editorial note.} All delivered mathematical text is the author's own, copied from the pinned Stacks Project chapter file algebra.tex; no statement or proof step was rewritten, reordered or removed.\par\medskip
\noindent\textit{Editorial note.} Preamble provenance, as required for this source: the frozen snapshot stages the chapter's own preamble file beside it (preamble.tex, 217 lines, 8116 bytes, the same shared Stacks-wide preamble the published chow and etale-cohomology cards use); the first line of the pinned chapter file is \textbackslash{}input\{preamble\}, so no preamble text lives inside the chapter itself. The card takes that file's lines 8-18 (the verbatim package and the three \textbackslash{}comment-based environments reference/slogan/history at lines 9-12, section 13.2, together with the Xy-pic setup) and its lines 154 to the runtime end of the file (the native theorem declarations and mathematical macros). Not one character was changed, no package was added, no macro was defined, and the excerpt serves only as preamble, never as a span coordinate.\par\medskip
\noindent\textit{Editorial note.} The same-chapter pointers are located at their \textbackslash{}label lines inside the pinned chapter file; the cross-chapter pointers are located in the project's own tag list staged in the snapshot (sources/stacks-chow/tags.txt), at the verbatim lines carrying their canonical tags.\par\medskip
\noindent\textit{Editorial note.} The counted proof refers to two same-chapter statements by label (the finiteness of the tensor product of finitely presented modules and the characterisation of Mittag-Leffler modules); both are recorded as navigation stubs pointed at their \textbackslash{}label lines, and the arguments consuming them stay inside the counted span.\par\medskip
\noindent\textit{Editorial note.} The counted spans contain no \textbackslash{}cite, so the card carries no bibliography.\par\medskip
\noindent\textit{Allowed standard prerequisites (editorial).}\begin{itemize}
\item The characterisation of Mittag-Leffler modules by the existence of compatible maps. Invoked twice in the counted proof (lines 21363-21369 and 21374-21375). Converts the Mittag-Leffler property of the two systems into the existence of the maps whose tensor product is then exhibited.
\item Finite presentation of the tensor product of finitely presented modules. Invoked in the counted proof at lines 21359-21361. Shows the approximating modules of the tensor product system are finitely presented.
\end{itemize}\medskip
\ifdefined\hypertarget\hypertarget{sourcepointer-1}{}\fi
\noindent\textit{Source reference.} Original lemma 05BS: sources/stacks-algebra/pinned-original-source.tex, source lines 2281--2281. Navigation pointer for \textbackslash{}ref\{lemma-tensor-finiteness\} used inside the counted span (line 21361); the argument that consumes it is retained in the span.\par\medskip
\ifdefined\hypertarget\hypertarget{sourcepointer-2}{}\fi
\noindent\textit{Source reference.} Original lemma 059E: sources/stacks-algebra/pinned-original-source.tex, source lines 21217--21217. Navigation pointer for \textbackslash{}ref\{proposition-ML-characterization\} used inside the counted span (line 21365); the argument that consumes it is retained in the span.\par\medskip
\begin{lemma}
\label{lemma-tensor-ML-modules}
If $R$ is a ring and $M$, $N$ are Mittag-Leffler modules over $R$,
then $M \otimes_R N$ is a Mittag-Leffler module.
\end{lemma}

\begin{proof}
Write $M = \colim_{i \in I} M_i$ and $N = \colim_{j \in J} N_j$
as directed colimits of finitely presented $R$-modules.
Denote $f_{ii'} : M_i \to M_{i'}$ and $g_{jj'} : N_j \to N_{j'}$ the
transition maps. Then $M_i \otimes_R N_j$ is a finitely presented
$R$-module (see
Lemma \ref{lemma-tensor-finiteness}),
and $M \otimes_R N = \colim_{(i, j) \in I \times J} M_i \otimes_R M_j$.
Pick $(i, j) \in I \times J$. By the definition of a Mittag-Leffler module
we have
Proposition \ref{proposition-ML-characterization} (3)
for both systems. In other words there exist $i' \geq i$ and $j' \geq j$
such that for every choice of $i'' \geq i$ and $j'' \geq j$ there exist
maps $a : M_{i''} \to M_{i'}$ and $b : M_{j''} \to M_{j'}$ such that
$f_{ii'} = a \circ f_{ii''}$ and $g_{jj'} = b \circ g_{jj''}$.
Then it is clear that
$a \otimes b : M_{i''} \otimes_R N_{j''} \to M_{i'} \otimes_R N_{j'}$
serves the same purpose for the system
$(M_i \otimes_R N_j, f_{ii'} \otimes g_{jj'})$.
Thus by the characterization
Proposition \ref{proposition-ML-characterization} (3)
we conclude that $M \otimes_R N$ is Mittag-Leffler.
\end{proof}

\end{document}
